\documentclass[main]{subfiles}

\begin{document}
% ref: claim_3_1 (proof)
% title: Untitled
% migrated from Section3_1/tex_proofs/claim_3_1_JNodeRestrictions.tex

% claim_3_1 — standalone TeX proof.
% This claim is stated in the LN as a \begin{Rem} without an accompanying
% \begin{proof} block; we construct the proof here by unfolding the
% definitions \ref{def-cdmg} and \ref{not-cdmg}.

\begin{claimmark}
\begin{Rem}
With the notations \ref{not-cdmg} the restrictions in definition \ref{def-cdmg} mean that the nodes $j \in J$ will not have any arrowheads pointing towards them: $ j \hus v \notin G$. Nodes $j \in J$ can only point towards nodes $v \in V$: edges $j \tuh v$ are allowed. Furthermore, no two nodes in $J$ are adjacent.
\end{Rem}
\end{claimmark}

\begin{proof}
The remark has three parts; each is obtained by unfolding the notation of
\ref{not-cdmg} against the type signatures of $E$ and $L$ given in
\ref{def-cdmg}. Recall from \ref{def-cdmg} that
\[
  E \ins (J \cup V) \times V
  \qquad\text{and}\qquad
  L \ins V \times V,
\]
so the second coordinate of any directed edge lies in $V$, and \emph{both}
coordinates of any bidirected edge lie in $V$. The sets $J$ and $V$ are
disjoint, hence no element of $J$ lies in $V$.

\medskip
\noindent\textbf{(1) No arrowhead points into a node $j \in J$.}
Fix $j \in J$ and $v \in G$. By \ref{not-cdmg}, $j \hus v \in G$ means
$j \hut v \in G$ or $j \huh v \in G$.
\begin{itemize}
  \item If $j \hut v \in G$, then by \ref{not-cdmg} $(v, j) \in E$. But by
    \ref{def-cdmg} the second coordinate of any element of $E$ lies in $V$,
    so $j \in V$, contradicting $j \in J$ and the disjointness $J \cap V = \emptyset$.
  \item If $j \huh v \in G$, then by \ref{not-cdmg} $(j, v) \in L$. But by
    \ref{def-cdmg} the first coordinate of any element of $L$ lies in $V$,
    so again $j \in V$, contradicting $j \in J$.
\end{itemize}
Both disjuncts are impossible, so $j \hus v \notin G$.

\medskip
\noindent\textbf{(2) Edges $j \tuh v$ with $j \in J$ and $v \in V$ are allowed.}
By \ref{not-cdmg}, $j \tuh v \in G$ means $(j,v) \in E$, and by
\ref{def-cdmg} $E \ins (J \cup V) \times V$. Since $j \in J \ins J \cup V$
and $v \in V$, the pair $(j,v)$ lives in the ambient set $(J \cup V) \times V$
that $E$ is a subset of; nothing in \ref{def-cdmg} excludes such a pair from
$E$. (This is an observation about which edges the definition \emph{permits},
not a theorem; it is recorded for symmetry with (1) and (3).)

\medskip
\noindent\textbf{(3) No two nodes in $J$ are adjacent.}
Fix $j_1, j_2 \in J$. By `def_3_3', $j_1$ and $j_2$ are adjacent in $G$ iff
$j_1 \sus j_2 \in G$, i.e.\ (by \ref{not-cdmg}) iff at least one of the three
relations
\[
  j_1 \tuh j_2 \in G, \qquad j_1 \hut j_2 \in G, \qquad j_1 \huh j_2 \in G
\]
holds. We rule out each in turn.
\begin{itemize}
  \item $j_1 \tuh j_2 \in G$ means $(j_1, j_2) \in E$. By \ref{def-cdmg} the
    second coordinate lies in $V$, so $j_2 \in V$, contradicting
    $j_2 \in J$ and $J \cap V = \emptyset$.
  \item $j_1 \hut j_2 \in G$ means $(j_2, j_1) \in E$. By the same reasoning
    applied to the second coordinate, $j_1 \in V$, contradiction.
  \item $j_1 \huh j_2 \in G$ means $(j_1, j_2) \in L$. By \ref{def-cdmg}
    $L \ins V \times V$, so $j_1 \in V$ (and $j_2 \in V$), contradiction.
\end{itemize}
None of the three disjuncts can hold, so $j_1 \sus j_2 \notin G$ and
therefore $j_1$ and $j_2$ are not adjacent. \qed
\end{proof}

\end{document}
